%--------------------------------------------------------------------------------
%
%
%--------------------------------------------------------------------------------
\unnumberedChapter{Twin-64 Simulator}

The Twin-64 simulator is an interactive program that runs a configured Twin-64 system. Running the simulator simply means launching the simulator application in a terminal window. During the execution of the configured system, data can be displayed and modified. The simulator features a processor model, system memory and a simulated I/O system. The simulator terminal screen displays system data and allows the monitoring, modifying and tracing of the system execution. The simulator environment is a set of configured modules. A module is a processor, memory or I/O module. The following figure shows a high level view of a configured system.

\begin{figure}[H]
\centering
\begin{tikzpicture}[x=1cm, y=1cm, scale=0.7, transform shape]

    % \draw[help lines, gray!70, dashed] (0,0) grid (20,10);
    
    \node[  tsRoundedRectangle, 
            minimum width=4cm,
            minimum height=8cm,
            line width = 1.3pt,  
            fill=gray!20,
            anchor=south west] at (1,1) { \LARGE Simulator \\ \LARGE Driver};
            
    \draw[->, gray, thick] (5, 7) -- (7, 7);
                
    \node[  tsRoundedRectangle, 
            minimum width=10cm,
            minimum height=8cm,  
             line width = 1.3pt, 
            fill=gray!20,
            anchor=south west] at (7,1) { };
            
    \draw[<-, gray, thick] (5, 3) -- (7, 3);
    
    \node[  tsRectangle, 
            minimum width=2cm,
            minimum height=2cm,  
            fill=white,
            anchor=south west] at (9.5,6) { Proc 1};
            
   	\node[  tsRectangle, 
            minimum width=2cm,
            minimum height=2cm,  
            fill=white,
            anchor=south west] at (12.5,6) { Proc 2};
            
    \draw[<->, gray, thick] (8, 5) -- ( 16, 5);
    
    \node[  tsRectangle, 
            minimum width=2cm,
            minimum height=2cm,  
            fill=white,
            anchor=south west] at (8, 2) { Memory};
            
    \node[  tsRectangle, 
            minimum width=2cm,
            minimum height=2cm,  
            fill=white,
            anchor=south west] at (11, 2) { I/O};
            
    \node[  tsRectangle, 
            minimum width=2cm,
            minimum height=2cm,  
            fill=white,
            anchor=south west] at (14, 2) { I/O};
            
    \draw[<->, black, thick] (10.5, 5) -- (10.5, 6);
    \draw[<->, black, thick] (13.5, 5) -- (13.5, 6);
    \draw[<->, black, thick] (9, 4) -- (9, 5);
    \draw[<->, black, thick] (12, 4) -- (12, 5);
	\draw[<->, black, thick] (15, 4) -- (15, 5);

\end{tikzpicture}
\caption{Simulator Environment}
\label{fig:simulator-environment}
\end{figure}

The figure above shows a two processor system with a memory module and two I/O modules an a shared \textbf{system bus}. The simulator driver issues commands to the simulation environment. Examples are modifying a memory cell content, or single stepping the processors. The simulated environment provides execution data to the driver. 

The data is generally shown in \textbf{simulator windows}. The main window is the command window. The simulator displays them arranged within a single terminal screen. Each window begins with a banner line shown in inverse video. The banner identifies the window and indicates its current state. The last window on the terminal screen is always the command window.

The simulator normally operates with multiple windows visible, but it can also be placed into \textbf{line mode}. In this mode, only the command window is displayed. Commands are entered at the prompt line, and all output is shown sequentially, one line at a time.
